% Options for packages loaded elsewhere
\PassOptionsToPackage{unicode}{hyperref}
\PassOptionsToPackage{hyphens}{url}
%
\documentclass[
  11pt,
]{article}
\usepackage{lmodern}
\usepackage{amssymb,amsmath}
\usepackage{ifxetex,ifluatex}
\ifnum 0\ifxetex 1\fi\ifluatex 1\fi=0 % if pdftex
  \usepackage[T1]{fontenc}
  \usepackage[utf8]{inputenc}
  \usepackage{textcomp} % provide euro and other symbols
\else % if luatex or xetex
  \usepackage{unicode-math}
  \defaultfontfeatures{Scale=MatchLowercase}
  \defaultfontfeatures[\rmfamily]{Ligatures=TeX,Scale=1}
\fi
% Use upquote if available, for straight quotes in verbatim environments
\IfFileExists{upquote.sty}{\usepackage{upquote}}{}
\IfFileExists{microtype.sty}{% use microtype if available
  \usepackage[]{microtype}
  \UseMicrotypeSet[protrusion]{basicmath} % disable protrusion for tt fonts
}{}
\makeatletter
\@ifundefined{KOMAClassName}{% if non-KOMA class
  \IfFileExists{parskip.sty}{%
    \usepackage{parskip}
  }{% else
    \setlength{\parindent}{0pt}
    \setlength{\parskip}{6pt plus 2pt minus 1pt}}
}{% if KOMA class
  \KOMAoptions{parskip=half}}
\makeatother
\usepackage{xcolor}
\IfFileExists{xurl.sty}{\usepackage{xurl}}{} % add URL line breaks if available
\IfFileExists{bookmark.sty}{\usepackage{bookmark}}{\usepackage{hyperref}}
\hypersetup{
  hidelinks,
  pdfcreator={LaTeX via pandoc}}
\urlstyle{same} % disable monospaced font for URLs
\usepackage[margin=1in]{geometry}
\setlength{\emergencystretch}{3em} % prevent overfull lines
\providecommand{\tightlist}{%
  \setlength{\itemsep}{0pt}\setlength{\parskip}{0pt}}
\setcounter{secnumdepth}{-\maxdimen} % remove section numbering
\usepackage{esint}
\setmathfont{Asana Math}
\setlength{\emergencystretch}{3em}

\author{}
\date{}

\begin{document}

\hypertarget{from-pointwise-ricci-control-to-finite-monodromy-and-exact-realization}{%
\section{From pointwise Ricci control to finite monodromy and exact
realization}\label{from-pointwise-ricci-control-to-finite-monodromy-and-exact-realization}}

24 September 2026. This is a composition theorem for a specified
geometric class, not a theorem that arbitrary near-Einstein
four-manifolds admit a surface fibration. Its global input is the
thick-fiber cycle theorem in
\href{20260924_cycle_holonomy_realization.md}{the companion note}. Its
curvature identity is classical and was already used in
\href{../mathematics/current/round20_construction.md}{Round 20, Sections
5--7}. The new project-level conclusion is a quantitative
\textbf{sufficient condition for exact existence}, including
zero-signature bundles: small pointwise Ricci error detects finite
mapping-class monodromy, instead of assuming it.

\hypertarget{statement}{%
\subsection{1. Statement}\label{statement}}

Let \[
        \Sigma_f\longrightarrow M\xrightarrow{\pi}B
\] be a smooth oriented surface bundle, with fiber genus \(f\ge2\) and
closed oriented base of genus at least two. A Riemannian metric \(g\)
makes \(\pi:(M,g)\to(B,b)\) a Riemannian submersion. Write \(h_y\) for
its vertical metric, \(\mathcal H\) for its orthogonal horizontal
distribution, and \(K_b\) for the Gaussian curvature of the base. Assume
the fibers are minimal. They need not be totally geodesic, intrinsically
hyperbolic, or mutually isometric. The horizontal distribution need not
be integrable.

Fix a reference fiber over \(y_0\). Uniformize its conformal class to
curvature \(-1\), and assume its hyperbolic systole is at least
\(\iota>0\). Define \[
\begin{split}
 \delta&=\|\operatorname{Ric}_g^{\sharp}+I\|_{L^\infty,\mathrm{op}},\\
 \eta&=\|(K_b+1)_+\|_{L^\infty(B)},\qquad
 D_B=\operatorname{diam}(B,b).
\end{split}                                                    \tag{1}
\] The positive part in \(\eta\) is intentional; no two-sided curvature
closeness of the base is needed for the estimate below.

\textbf{Theorem 1 (Ricci-to-monodromy realization).} There is
\(\epsilon_{f,\iota}>0\), depending only on fiber genus and the
reference systole, such that \[
                 2D_B\sqrt{\delta+\eta}<\epsilon_{f,\iota}       \tag{2}
\] implies that the mapping-class monodromy of the bundle has finite
image. Consequently the \textbf{same oriented smooth manifold} \(M\)
carries a negative Kähler--Einstein metric \(\widehat g\), with \[
 \operatorname{Ric}_{\widehat g}=-\widehat g,
 \qquad (\widetilde M,\widetilde{\widehat g})
          \cong\mathbb H^2\times\mathbb H^2.                  \tag{3}
\] The theorem imposes no bound on total-space diameter, total volume,
injectivity radius, full curvature, or a spectral inverse. It does not
assert closeness of \(\widehat g\) to \(g\).

In particular, when \(K_b=-1\), condition (2) is simply \[
          \delta<\frac{\epsilon_{f,\iota}^2}{4D_B^2}.          \tag{4}
\] Thus, within this explicitly specified input class, the output is
already the Einstein existence requested by the negative side of Yau26.
The additional fibration, minimality, thickness and base-scale
hypotheses are the restriction; obtaining KE rather than a general
Einstein metric is not a defect of the existence conclusion.

\hypertarget{horizontal-ricci-controls-conformal-deformation}{%
\subsection{2. Horizontal Ricci controls conformal
deformation}\label{horizontal-ricci-controls-conformal-deformation}}

For a horizontal vector \(X\), define the symmetric vertical
endomorphism \(S_X\) by \[
            h(S_XU,V)=\langle(\nabla_UV)^{\mathcal H},X\rangle.
\] Let \(A_XY=\tfrac12[X,Y]^{\mathcal V}\) for basic horizontal vector
fields. Minimality is exactly \(\operatorname{tr}S_X=0\) for every
horizontal \(X\). For a unit horizontal vector, O'Neill's contracted
curvature formulas give \[
 \operatorname{Ric}_g(X,X)
       =K_b-|S_X|_{\mathrm{HS}}^2
                -2\sum_a|A_Xe_a|^2.                         \tag{5}
\] Here \(e_a\) is a horizontal orthonormal basis. The derivative of the
fiber mean-curvature vector in the general formula is zero because that
vector is identically zero, not because its integral vanishes. To check
the coefficients, horizontal sectional curvatures contribute
\(K_b-3\sum_a|A_Xe_a|^2\); the mixed sectional curvatures contribute
\(\sum_a|A_Xe_a|^2-|S_X|_{\mathrm{HS}}^2\). This is the same identity
used in Round 20; our sign choice for \(S\) has no effect on (5). The
primary source is O'Neill,
\href{https://doi.org/10.1307/mmj/1028999604}{\emph{The fundamental
equations of a submersion}}.

Since the horizontal compression of \(\operatorname{Ric}_g+g\) has
operator norm at most \(\delta\), (5) yields, pointwise, \[
\begin{split}
 |S_X|_{\mathrm{HS}}^2+2\sum_a|A_Xe_a|^2
   &=(K_b+1)-(\operatorname{Ric}_g+g)(X,X)\\
   &\le\eta+\delta.
\end{split}                                                \tag{6}
\] In particular, for an arbitrary horizontal vector, \[
                 |S_X|_{\mathrm{op}}
           \le\sqrt{\delta+\eta}\,|X|.                    \tag{7}
\] No vertical or mixed Ricci equation is discarded in the hypothesis:
the full Ricci defect in (1) implies the horizontal bound. Conversely we
make no claim that the horizontal bound alone makes an arbitrary trial
metric near-Einstein. Intrinsic fiber curvature is absent from (5),
which is why no fiber uniformization hypothesis is required away from
\(y_0\).

\hypertarget{integrating-the-actual-deformation-around-a-loop}{%
\subsection{3. Integrating the actual deformation around a
loop}\label{integrating-the-actual-deformation-around-a-loop}}

Let \(\gamma:[0,1]\to B\) be a piecewise smooth loop based at \(y_0\).
Its horizontal transport \(P_t:F_{y_0}\to F_{\gamma(t)}\) is defined for
the entire interval because the bundle is compact. Set
\(k_t=P_t^*h_{\gamma(t)}\). Differentiating on the fixed marked fiber,
\[
                 k_t^{-1}\dot k_t=-2P_t^*S_{\dot\gamma^{\mathcal H}}.
                                                                  \tag{8}
\] The trace is zero by minimality. Write \(x_t\) for the corresponding
point in Teichmüller space, with metric normalized by
\(d_T=\tfrac12\inf\log K\).

For clarity the deformation estimate does not require differentiating
the uniformization map. For any nonzero tangent vector on the reference
fiber, \[
 \left|\frac{d}{dt}\log |v|_{k_t}\right|
       \le\sup_{F_{\gamma(t)}}|S_{\dot\gamma^{\mathcal H}}|_{\mathrm{op}}.
\] After integration, the ratio between the largest and smallest
singular stretch of the identity map is at most
\(\exp(2\int\sup|S|\,dt)\). This ratio is its quasiconformal dilatation.
It follows that \[
\begin{split}
 d_T(x_0,x_1)
 &\le\int_0^1\sup_{F_{\gamma(t)}}
            |S_{\dot\gamma^{\mathcal H}}|_{\mathrm{op}}\,dt\\
 &\le\sqrt{\delta+\eta}\,\operatorname{Length}_b(\gamma).
\end{split}                                                     \tag{9}
\] The endpoint is the action on \(x_0\) of the monodromy class, or its
inverse according to the marking convention. Distances are the same.
Connection holonomy itself need not be a homomorphism into
diffeomorphisms when the connection is curved; its
\textbf{mapping-class} holonomy is the topological monodromy
homomorphism. Formula (9) controls that homomorphism without falsely
imposing flatness of the connection.

\hypertarget{short-generators-and-a-common-finite-group}{%
\subsection{4. Short generators and a common finite
group}\label{short-generators-and-a-common-finite-group}}

We use the following elementary global fact.

\textbf{Lemma 2.} For a closed connected Riemannian manifold \(B\), any
basepoint \(y_0\), and \(a>0\), its fundamental group has a finite
generating set represented by loops of length less than
\(2\operatorname{diam}(B)+a\).

\textbf{Proof.} Subdivide any based loop into arcs of length less than
\(a\). Join each subdivision vertex to \(y_0\) by a minimizing path of
length at most \(D_B\). The original loop is the product of the
resulting triangular based loops; cancellation is literal along the
joining paths. Each triangular loop has length less than \(2D_B+a\).
Only finitely many based homotopy classes admit a representative with
that length bound: lifting to the complete universal cover associates
each class to a distinct deck translate of a fixed lift of \(y_0\) in a
compact ball. The deck orbit is discrete and the action is proper. Thus
the collection of all these short classes is a finite generating set.
\(\square\)

The companion theorem supplies the following precise Teichmüller fact:
for every \(f,\iota\), there exists \(\epsilon_{f,\iota}>0\) such that,
at every marked hyperbolic surface of systole at least \(\iota\), the
subgroup generated by \textbf{all} mapping classes moving it less than
\(\epsilon_{f,\iota}\) is finite. This is stronger than saying that each
such class separately has finite order. It follows from proper
discontinuity of the mapping-class action and compactness of the thick
moduli quotient, using neighborhoods of a \textbf{common} finite
stabilizer. The companion note proves the uniform statement and treats
the possible ineffective finite kernel in genus two.

By the strict inequality (2), choose \(a>0\) sufficiently small that \[
              (2D_B+a)\sqrt{\delta+\eta}<\epsilon_{f,\iota}.
\] Apply (9) to the finite generating set of Lemma 2. All its monodromy
elements lie in the same finite subgroup supplied by the thick-point
lemma. Therefore the entire monodromy image is finite. This is the
global step that a pointwise estimate for \(S\) alone does not provide.

\hypertarget{realization-on-the-original-smooth-bundle}{%
\subsection{5. Realization on the original smooth
bundle}\label{realization-on-the-original-smooth-bundle}}

Let \(G\) be the finite monodromy image. Nielsen realization gives a
hyperbolic metric \(h_F\) on the model fiber and an action
\(G\hookrightarrow\operatorname{Isom}^+(\Sigma_f,h_F)\) realizing those
mapping classes. Let \(\widehat B\to B\) be the connected regular cover
corresponding to the kernel of monodromy. Its deck group is \(G\).
Uniformize any conformal structure on the oriented base to a hyperbolic
metric \(\widehat b\) and lift it to \(\widehat B\). Then \[
          (\widehat B\times\Sigma_f)/G                     \tag{10}
\] is a smooth diagonal quotient; the action is free because it is free
on the base factor. It carries the quotient of the product metric, with
Ricci tensor equal to minus the metric, and the quotient complex
structure is Kähler.

The identity component of the orientation-preserving diffeomorphism
group of a closed surface of genus at least two is contractible. Thus
oriented smooth surface bundles are classified by their mapping-class
monodromy; there is no additional identity-component twisting over
\(B\). The quotient bundle (10) and the original bundle are smoothly
isomorphic, since their monodromy representations coincide. Pulling back
the quotient metric proves (3) on the original \(M\). The companion note
gives the classical Nielsen and surface-diffeomorphism references.
Neither the original horizontal distribution nor its full diffeomorphism
holonomy has to agree with the resulting product quotient.

\hypertarget{a-quantified-escape-alternative}{%
\subsection{6. A quantified escape
alternative}\label{a-quantified-escape-alternative}}

Taking the contrapositive gives a useful theorem about actual inputs,
not an assertion of general Einstein nonexistence.

\textbf{Corollary 3.} If the mapping-class monodromy image is infinite
and the reference fiber is \(\iota\)-thick, then every
minimal-submersion metric as above satisfies \[
          (\delta+\eta)D_B^2\ge\frac{\epsilon_{f,\iota}^2}{4}.
                                                               \tag{11}
\] For hyperbolic bases and \(\delta_j\to0\), a sequence of such metrics
with fixed fiber genus and a uniformly thick reference fiber therefore
has \[
          D_{B,j}\ge\frac{\epsilon_{f,\iota}}{2\sqrt{\delta_j}}
                   \longrightarrow\infty.                  \tag{12}
\] Since a Riemannian submersion is distance nonincreasing,
\(\operatorname{diam}(M,g)\ge D_B\). Alternatively, with bounded
\(D_B\), infinite monodromy and vanishing defect force the reference
conformal systole to approach zero along any subsequence for which it
could otherwise have a positive lower bound.

These alternatives distinguish two real escape mechanisms: long
horizontal transport and thin fiber conformal geometry. The companion
note supplies Dehn-twist examples separating them. They do not say that
the same smooth manifold has no Einstein metric outside this submersion
class, nor that every near-Einstein manifold has either of these two
geometric descriptions.

\hypertarget{what-changes-and-what-does-not}{%
\subsection{7. What changes, and what does
not}\label{what-changes-and-what-does-not}}

Round 20 already proved a signature-energy obstruction for minimal
hyperbolic-fiber submersions and recovered the classical exact adapted
Einstein classification. For a signature-zero bundle its signature lower
bound vanishes. The present argument instead uses the actual
mapping-class action and one thick conformal fiber. It can detect
infinite monodromy even when signature is zero; more importantly, it
proves existence when the smallness threshold is met.

This is not a new construction of bidisk quotients: those are classical.
It is a geometric criterion that \textbf{recognizes} their smooth bundle
types from the input metric, with finite monodromy deduced rather than
assumed. No assertion of worldwide priority is made. The first-line
connection is that local or average Kähler-type structure requires
global transport information to produce existence. Here that information
is obtained from a definite horizontal Ricci identity and a thick-point
monodromy threshold, rather than from another local calibration
coefficient.

The remaining step toward arbitrary Yau26 inputs is substantial: neither
an actual surface submersion nor fiber minimality nor thickness is
supplied by the general pointwise Ricci hypothesis. The existing
collapsed-bulk theorem does not automatically supply them across its
exceptional set. This theorem completes the existence argument for the
stated input class without claiming that missing classification.

\end{document}
